\documentclass[10pt,letterpaper]{article}
\usepackage{amsmath, amssymb}
\usepackage[margin=0.85in]{geometry}
\usepackage{multicol}
\usepackage{enumitem}
\usepackage{titlesec}
\usepackage{parskip}

\titleformat{\section}{\large\bfseries}{}{0em}{}[\titlerule]
\titleformat{\subsection}{\normalsize\bfseries}{}{0em}{}

\setlength{\columnsep}{1.2em}
\setlist[enumerate]{leftmargin=*, label=\textbf{\arabic*.}}

\begin{document}

\begin{center}
  {\LARGE\bfseries ICTM Regional Division A --- Algebra II}\\[4pt]
  {\large Formula \& Reference Sheet}
\end{center}

\vspace{0.5em}

% ─────────────────────────────────────────────
\section{Complex Numbers}

\begin{enumerate}

\item \textbf{Imaginary unit and powers of $i$}

$$i = \sqrt{-1}, \quad i^2 = -1, \quad i^3 = -i, \quad i^4 = 1$$

The powers cycle with period 4. To find $i^n$, compute $n \bmod 4$:
$$i^n = i^{n \bmod 4}$$

\item \textbf{Reducing large powers} (e.g.\ $i^{2026}$)

Divide the exponent by 4 and use the remainder:
$$2026 = 4 \times 506 + 2 \implies i^{2026} = i^2 = -1$$

\item \textbf{Sum of consecutive powers}

Any four consecutive integer powers of $i$ sum to 0. Use this to collapse large sums:
$$i^2 + i^3 + i^4 + i^5 + i^6 + i^7 + i^8 = (i^2+i^3+i^4+i^5)+(i^6+i^7+i^8) = 0 + i^6+i^7+i^8$$

\item \textbf{Operations with complex numbers}

$$(a+bi)+(c+di) = (a+c)+(b+d)i$$
$$(a+bi)(c+di) = (ac-bd)+(ad+bc)i$$

\item \textbf{Modulus (absolute value)}

$$|a+bi| = \sqrt{a^2+b^2}$$

\end{enumerate}

% ─────────────────────────────────────────────
\section{Logarithms \& Exponents}

\begin{enumerate}

\item \textbf{Definition}

$$\log_b x = y \iff b^y = x, \quad b > 0,\; b \ne 1,\; x > 0$$

\item \textbf{Core identities}

$$\log_b 1 = 0 \qquad \log_b b = 1 \qquad b^{\log_b x} = x$$

\item \textbf{Laws of logarithms}

$$\log_b(MN) = \log_b M + \log_b N$$
$$\log_b\!\left(\frac{M}{N}\right) = \log_b M - \log_b N$$
$$\log_b(M^p) = p\log_b M$$

\item \textbf{Change of base}

$$\log_b x = \frac{\log x}{\log b} = \frac{\ln x}{\ln b}$$

\item \textbf{Common and natural logs}

$$\log x \equiv \log_{10} x \qquad \ln x \equiv \log_e x$$
$$\ln e^x = x \qquad e^{\ln x} = x$$

\item \textbf{Solving exponential equations}

Equal bases $\Rightarrow$ set exponents equal. Otherwise take $\log$ of both sides and apply the power rule.

\item \textbf{Remainder when $10^{20}$ is divided by 6}

Note $10^n \equiv 4 \pmod{6}$ for all $n \ge 1$, so the remainder is \textbf{4}.

\end{enumerate}

% ─────────────────────────────────────────────
\section{Sequences \& Series}

\begin{enumerate}

\item \textbf{Arithmetic sequence}

$$a_n = a_1 + (n-1)d, \qquad d = a_2 - a_1$$

\item \textbf{Sum of $n$ terms (arithmetic)}

$$S_n = \frac{n}{2}(a_1 + a_n) = \frac{n}{2}\bigl(2a_1+(n-1)d\bigr)$$

\item \textbf{Geometric sequence}

$$a_n = a_1 \cdot r^{\,n-1}, \qquad r = \frac{a_2}{a_1}$$

\item \textbf{Sum of $n$ terms (geometric, $r \ne 1$)}

$$S_n = \frac{a_1(1-r^n)}{1-r}$$

\item \textbf{Infinite geometric series ($|r| < 1$)}

$$S = \frac{a_1}{1-r}$$

\item \textbf{Arithmetic mean of two values}

$$\text{AM} = \frac{a+b}{2}$$

\item \textbf{Geometric mean of two values}

$$\text{GM} = \sqrt{ab}, \qquad a,b > 0$$

\item \textbf{Harmonic mean of two values}

$$\text{HM} = \frac{2ab}{a+b}$$

\item \textbf{AM--GM inequality}

$$\frac{a+b}{2} \ge \sqrt{ab}, \qquad a,b \ge 0$$
Equality holds when $a = b$.

\item \textbf{Sum of even integers $2 + 4 + \cdots + 2n$}

$$S = n(n+1)$$

\item \textbf{Recursive sequences}

Always compute term-by-term using the given rule, e.g.\ $a_{n+1} = a_n^{\,2} - 2a_n$ with $a_1 = 2$.

\item \textbf{Pell numbers} ($P_0=0,\; P_1=1$)

$$P_n = 2P_{n-1} + P_{n-2}, \qquad n \ge 2$$

First several: $0,\,1,\,2,\,5,\,12,\,29,\,70,\ldots$

\end{enumerate}

% ─────────────────────────────────────────────
\section{Polynomials \& Roots}

\begin{enumerate}

\item \textbf{Remainder theorem}

When $f(x)$ is divided by $(x-c)$, the remainder equals $f(c)$.

\item \textbf{Factor theorem}

$(x-c)$ is a factor of $f(x) \iff f(c) = 0$.

\item \textbf{Quadratic formula}

$$x = \frac{-b \pm \sqrt{b^2-4ac}}{2a}$$

\item \textbf{Discriminant $\Delta = b^2 - 4ac$}

\begin{itemize}
  \item $\Delta > 0$: two distinct real roots
  \item $\Delta = 0$: one real (double) root
  \item $\Delta < 0$: two complex conjugate roots
\end{itemize}

\item \textbf{Vieta's formulas (quadratic $ax^2+bx+c=0$, roots $r_1, r_2$)}

$$r_1 + r_2 = -\frac{b}{a} \qquad r_1\,r_2 = \frac{c}{a}$$

\item \textbf{Sum of squares of roots (quadratic)}

$$r_1^2 + r_2^2 = (r_1+r_2)^2 - 2r_1 r_2$$

\item \textbf{Difference of roots (quadratic)}

$$|r_1 - r_2| = \frac{\sqrt{b^2-4ac}}{|a|}$$

\item \textbf{Tangent to the $x$-axis condition}

A parabola $y = ax^2+bx+c$ is tangent to the $x$-axis $\iff \Delta = 0 \iff b^2 = 4ac$.

\item \textbf{Finding $k$ or $w$ when two polynomials share a common root}

Subtract the two equations to eliminate $x^3$ and higher-degree shared terms, then solve for the shared root, and substitute back.

\item \textbf{Polynomial with given zeros}

If zeros are $r_1, r_2, \ldots, r_n$ and leading coefficient is 1:
$$f(x) = (x-r_1)(x-r_2)\cdots(x-r_n)$$

\end{enumerate}

% ─────────────────────────────────────────────
\section{Combinatorics \& Probability}

\begin{enumerate}

\item \textbf{Permutations} (ordered, $r$ chosen from $n$)

$$P(n,r) = \frac{n!}{(n-r)!}$$

\item \textbf{Combinations} (unordered, $r$ chosen from $n$)

$$C(n,r) = \binom{n}{r} = \frac{n!}{r!\,(n-r)!}$$

\item \textbf{Arrangements with repeated elements}

$$\frac{n!}{n_1!\,n_2!\,\cdots\,n_k!}$$

\item \textbf{Binomial theorem}

$$(x+y)^n = \sum_{k=0}^{n} \binom{n}{k} x^{n-k}\,y^{\,k}$$

General term ($(k+1)$-th term):
$$T_{k+1} = \binom{n}{k} x^{n-k}\,y^{\,k}$$

\item \textbf{Sum of binomial coefficients}

$$\sum_{k=0}^{n}\binom{n}{k} = 2^n$$

This is the sum of \emph{numerical} coefficients when $x=y=1$.

\item \textbf{Complement rule}

$$\Pr(A^c) = 1 - \Pr(A)$$

\item \textbf{Probability of at least one success in $n$ independent trials}

$$\Pr(\text{at least one}) = 1-(1-p)^n$$

\item \textbf{Inclusion--exclusion (two sets)}

$$|A \cup B| = |A| + |B| - |A \cap B|$$

\item \textbf{Distributing $n$ distinct objects into $k$ distinct containers (equally likely)}

Total outcomes: $k^n$. For equal distribution ($n/k$ per container):
$$\text{Favorable} = \frac{n!}{\left(\frac{n}{k}\right)!^k}$$

\end{enumerate}

% ─────────────────────────────────────────────
\section{Matrices \& Determinants}

\begin{enumerate}

\item \textbf{Matrix notation}

An $m \times n$ matrix has $m$ rows and $n$ columns:
$$A = \begin{bmatrix} a_{11} & a_{12} \\ a_{21} & a_{22} \end{bmatrix}$$

\item \textbf{$2\times2$ determinant}

$$\det\begin{bmatrix}a & b \\ c & d\end{bmatrix} = ad - bc$$

\item \textbf{$3\times3$ determinant (cofactor expansion along row 1)}

$$\det\begin{vmatrix}a & b & c \\ d & e & f \\ g & h & i\end{vmatrix}
= a(ei-fh) - b(di-fg) + c(dh-eg)$$

\item \textbf{$2\times2$ inverse}

$$A^{-1} = \frac{1}{\det A}\begin{bmatrix}d & -b \\ -c & a\end{bmatrix}$$

Exists only when $\det A \ne 0$.

\item \textbf{Identity matrix}

$$I_2 = \begin{pmatrix}1 & 0 \\ 0 & 1\end{pmatrix} \qquad I_3 = \begin{pmatrix}1 & 0 & 0 \\ 0 & 1 & 0 \\ 0 & 0 & 1\end{pmatrix}$$

$$A \cdot I = I \cdot A = A$$

\item \textbf{Matrix equation}

$$AX = B \implies X = A^{-1}B \quad (\text{when } \det A \ne 0)$$

Example:
$$\begin{bmatrix}7 & 6 \\ 4 & 2\end{bmatrix}\begin{bmatrix}k \\ w\end{bmatrix} = \begin{bmatrix}30 \\ 15\end{bmatrix}$$

\item \textbf{Augmented matrix for a linear system}

For $ax+by=e$, $cx+dy=f$:
$$\left[\begin{array}{cc|c} a & b & e \\ c & d & f \end{array}\right]$$

Solve by row reduction (elementary row operations: swap rows, scale a row, add a multiple of one row to another).

\item \textbf{Singular (non-invertible) matrix}

$$\det A = 0 \implies A \text{ is singular; the system } AX = B \text{ has no unique solution}$$

\item \textbf{Matrix multiplication} ($A$ is $m\times n$, $B$ is $n\times p$, result is $m\times p$)

$$(AB)_{ij} = \sum_{k=1}^{n} A_{ik}\,B_{kj}$$

Note: $AB \ne BA$ in general.

\item \textbf{System of equations as a matrix}

$$\begin{Bmatrix} 2x + y = 5 \\ x - 3y = -4 \end{Bmatrix} \longrightarrow \begin{bmatrix}2 & 1 \\ 1 & -3\end{bmatrix}\begin{bmatrix}x \\ y\end{bmatrix} = \begin{bmatrix}5 \\ -4\end{bmatrix}$$

\end{enumerate}

% ─────────────────────────────────────────────
\section{Conic Sections}

\begin{enumerate}

\item \textbf{Circle}

$$(x-h)^2 + (y-k)^2 = r^2, \quad \text{center } (h,k), \text{ radius } r$$

\item \textbf{Completing the square}

$$x^2 + bx = \left(x + \frac{b}{2}\right)^2 - \frac{b^2}{4}$$

\item \textbf{Ellipse (center $(h,k)$, $a > b > 0$)}

$$\frac{(x-h)^2}{a^2} + \frac{(y-k)^2}{b^2} = 1$$

Major axis length $= 2a$, minor axis length $= 2b$.

\item \textbf{Parabola opening right/left (vertex $(h,k)$)}

$$x - h = \frac{1}{4p}(y-k)^2$$

Focus at $(h+p,\, k)$, directrix $x = h-p$.

\item \textbf{Parabola opening up/down (vertex $(h,k)$)}

$$y - k = \frac{1}{4p}(x-h)^2$$

Focus at $(h,\, k+p)$, directrix $y = k-p$.

\item \textbf{Parabola with focus $(f_x, f_y)$ and directrix $y = d$}

The vertex is halfway between focus and directrix. The equation:
$$\frac{1}{4p}(x-h)^2 = y - k, \qquad p = f_y - k$$

\end{enumerate}

% ─────────────────────────────────────────────
\section{Functions \& Transformations}

\begin{enumerate}

\item \textbf{Transformation rule}

If $y = f(x)$ has a known point $(a, b)$, then $g(x) = A\,f\!\left(B(x-h)\right)+k$ maps it to:
$$\left(a \cdot \frac{1}{B} + h,\; A \cdot b + k\right)$$

Specifically: horizontal shift $+h$, horizontal scale $\frac{1}{B}$, vertical scale $A$, vertical shift $+k$.

\item \textbf{Vertex of $y = ax^2+bx+c$}

$$x_v = -\frac{b}{2a}, \qquad y_v = c - \frac{b^2}{4a}$$

\item \textbf{Maximum or minimum value (quadratic)}

$$y_{\max/\min} = \frac{4ac - b^2}{4a}$$

Maximum when $a < 0$; minimum when $a > 0$.

\item \textbf{Inverse function}

$f^{-1}$ exists iff $f$ is one-to-one. To find $f^{-1}(x)$: swap $x$ and $y$, then solve for $y$.
$$f(f^{-1}(x)) = x \qquad f^{-1}(f(x)) = x$$

\item \textbf{Composition}

$$(f \circ g)(x) = f(g(x))$$

\item \textbf{Rational function with $x$-intercept at $x = c$}

Set the numerator equal to 0 (and confirm $c$ is not also a zero of the denominator).

\end{enumerate}

% ─────────────────────────────────────────────
\section{Variation}

\begin{enumerate}

\item \textbf{Direct variation} ($y$ varies directly as $x$)

$$y = kx$$

\item \textbf{Inverse variation} ($y$ varies inversely as $x$)

$$y = \frac{k}{x} \qquad \Leftrightarrow \qquad xy = k$$

\item \textbf{Joint variation} ($z$ varies jointly as $x$ and $y$)

$$z = kxy$$

\item \textbf{Combined variation} (example: $x$ varies directly as $z$, inversely as $y$, jointly as $w^3$)

$$x = \frac{k\,z\,w^3}{y}$$

\item \textbf{Work rate problems}

$$\text{Rate} \times \text{Time} = \text{Work}$$

If $n$ workers complete $W$ units in $T$ hours: rate per worker $= \dfrac{W}{nT}$.

$$\text{Workers} \times \text{Rate} \times \text{Time} = \text{Work}$$

\end{enumerate}

% ─────────────────────────────────────────────
\section{Distance, Geometry \& Radicals}

\begin{enumerate}

\item \textbf{Distance between two points $(x_1,y_1)$ and $(x_2,y_2)$}

$$d = \sqrt{(x_2-x_1)^2 + (y_2-y_1)^2}$$

\item \textbf{Midpoint formula}

$$M = \left(\frac{x_1+x_2}{2},\; \frac{y_1+y_2}{2}\right)$$

\item \textbf{Distance from a point $(x_0,y_0)$ to a line $ax+by+c=0$}

$$d = \frac{|ax_0 + by_0 + c|}{\sqrt{a^2+b^2}}$$

\item \textbf{Distance between two parallel lines $ax+by+c_1=0$ and $ax+by+c_2=0$}

$$d = \frac{|c_1-c_2|}{\sqrt{a^2+b^2}}$$

\item \textbf{Space diagonal of a rectangular box ($l \times w \times h$)}

$$d = \sqrt{l^2 + w^2 + h^2}$$

\item \textbf{Simplifying radicals}

$$\sqrt{a} \cdot \sqrt{b} = \sqrt{ab} \qquad \frac{\sqrt{a}}{\sqrt{b}} = \sqrt{\frac{a}{b}}$$
$$\sqrt[n]{a} \cdot \sqrt[n]{b} = \sqrt[n]{ab}$$

\item \textbf{Rationalising the denominator}

$$\frac{k}{\sqrt{a}+\sqrt{b}} = \frac{k(\sqrt{a}-\sqrt{b})}{a-b}$$

\item \textbf{AM--GM (applied to optimisation)}

For $x + y = c$ (constant), the product $xy$ is maximised when $x = y = \dfrac{c}{2}$:
$$xy_{\max} = \frac{c^2}{4}$$

\end{enumerate}

\end{document}
